\section{术语与指标附录：把系统讨论落到可执行口径}

\subsection{关键术语对照}
为避免跨团队沟通歧义，这里把课程中涉及的核心概念统一为工程口径：
\begin{table}[H]
\centering
\begin{tabular}{p{2.6cm}p{4.0cm}p{4.4cm}}
\toprule
\textbf{术语} & \textbf{课程语境} & \textbf{工程落地定义} \\
\midrule
Agentic AI & 模型具备多步执行与工具调用能力 & 具状态、可规划、可执行、可恢复的任务系统 \\
RAG & 外部知识增强生成 & 可追溯检索+证据注入+答案生成管线 \\
Data Locality & 数据与算力的空间邻近性 & 任务调度时的数据/网络拓扑约束 \\
Exploration & 冲刺模型能力上限 & 容忍高成本的试验型研发阶段 \\
Exploitation & 追求部署效率与收益 & 以单位成本和稳定性为主导的生产阶段 \\
\bottomrule
\end{tabular}
\caption{课程高频术语的工程化定义}
\end{table}

\subsection{建议的周报指标模板}
如果团队已经在做 Agent 平台迭代，推荐每周至少跟踪以下指标，并按 ``质量、效率、稳定性'' 三类汇报：
\begin{longtable}{p{2.8cm}p{3.8cm}p{2.7cm}p{3.0cm}}
\toprule
\textbf{指标类别} & \textbf{指标} & \textbf{建议频率} & \textbf{目标趋势} \\
\midrule
质量 & 任务成功率、事实一致率、误操作率 & 日/周 & 成功率上升，误操作下降 \\
效率 & 平均成本/任务、GPU 利用率、缓存命中率 & 日/周 & 成本下降，利用率上升 \\
稳定性 & P95 延迟、错误峰值、MTTR & 日/周 & 尾延迟与 MTTR 下降 \\
治理 & 可审计覆盖率、权限违规次数 & 周/月 & 审计覆盖率上升，违规归零 \\
\bottomrule
\end{longtable}

\begin{knowledgebox}{为什么要把指标写进笔记}
课程知识只有转化为团队的例行观测和决策机制，才会变成真实能力。指标模板的作用是让 ``系统设计'' 进入持续迭代闭环，而不是停留在一次性讨论。
\end{knowledgebox}

\subsection{本章小结}
统一术语和指标口径，是跨研究、平台、产品团队协作的最低成本手段。它直接决定系统优化是否可持续、可复盘、可扩展。

